% SPDX-License-Identifier: CC-BY-SA-4.0
% Part: Lexical Analysis
% Section: From regular expressions to automata
% Exercise: Revision exercise

\begingroup

\begin{exercice}[Révisions -- Génération d'analyseur lexical]\label{exo:lexing/kleene/revision}

  Dessinez l'automate fini non déterministe obtenu en appliquant l'algorithme de Thompson à chacune des expressions rationnelles suivantes.
  \begin{question}
  \item $a^\star ba \mid b^\star aba$
  \item $(a \mid b) (a \mid bb^\star a)^\star bb^\star$
  \end{question}

  \begin{question}
  \item Traduire chaque automate ci-dessous en une expression rationnelle sur l'alphabet $\{a,b\}$.
  \end{question}
  
  \begin{tabular}{cccc}
    \begin{tikzpicture}[automaton, size=13mm, baseline=(1.base)]
      \state[initial  ] (1) at (0,0) {$1$}; 
      \state[         ] (2) at (1,0) {$2$}; 
      \state[accepting] (3) at (2,0) {$3$}; 
      \path (1) edge node {$a, b$} (2);
      \path (2) edge node {$a$}    (3);
    \end{tikzpicture}
    &
    \begin{tikzpicture}[automaton, size=13mm, baseline=(1.base)]
      \state[initial, accepting] (1) at (0,0) {$1$}; 
      \state                     (2) at (1,0) {$2$}; 
      \path (1) edge             node {$a$} (2);
      \path (2) edge[loop above] node {$a$} (2);
    \end{tikzpicture}
    &
    \begin{tikzpicture}[automaton, size=13mm, baseline=(1.base)]
      \state[initial, accepting] (1) at (0,0) {$1$}; 
      \state                     (2) at (1,0) {$2$}; 
      \path (1) edge[bend left]  node {$a$} (2);
      \path (2) edge[bend left]  node {$b$} (1);
      \path (1) edge[loop above] node {$b$} (1);
    \end{tikzpicture}
    &
    \begin{tikzpicture}[automaton, size=13mm, baseline=(1.base)]
      \state[initial]   (1) at (0,0) {$1$}; 
      \state            (2) at (1,0) {$2$}; 
      \state[accepting] (3) at (2,0) {$3$}; 
      \path (1) edge             node {$a$}    (2);
      \path (2) edge             node {$b$}    (3);
      \path (1) edge[loop above] node {$a, b$} (1);
      \path (3) edge[loop above] node {$a, b$} (3);
    \end{tikzpicture}
    \\
    \begin{tikzpicture}[automaton, size=13mm, baseline=(1.base)]
      \state[initial, accepting] (1) at (0,1) {$1$}; 
      \state                     (2) at (1,1) {$2$}; 
      \state                     (3) at (1,0) {$3$}; 
      \state[accepting]          (4) at (2,1) {$4$}; 

      \path (1) edge[loop above] node       {$b$} (1);
      \path (2) edge[loop above] node       {$a$} (2);
      \path (1) edge             node       {$a$} (2);
      \path (2) edge             node[swap] {$b$} (3);
      \path (3) edge             node       {$b$} (1);
      \path (4) edge             node[swap] {$a$} (2);
      \path (3) edge             node       {$a$} (4);
      \path (4) edge[bend left]  node       {$b$} (3);
    \end{tikzpicture}
    &
    \begin{tikzpicture}[automaton, size=13mm, baseline=(1.base)]
      \state[initial, accepting] (1) at (0,1) {$1$}; 
      \state                     (2) at (1,1) {$2$}; 
      \state                     (3) at (1,0) {$3$}; 
      \state[accepting]          (4) at (0,0) {$4$}; 

      \path (2) edge[loop above] node {$b$}   (2);
      \path (3) edge[loop below] node {$b$}   (3);
      \path (4) edge[loop left ] node {$b$}   (4);
      \path (2) edge             node {$a$}   (3);
      \path (3) edge             node {$a$}   (4);
      \path (4) edge             node {$a,b$} (1);
      \path (3) edge             node {$b$}   (1);
      \path (1) edge[bend left]  node {$a$}   (2);
      \path (2) edge             node {$b$}   (1);
    \end{tikzpicture}
    &
    \begin{tikzpicture}[automaton, size=13mm, baseline=(2.base)]
      \state[initial, accepting] (1) at ( 90:12mm) {$1$}; 
      \state                     (2) at (162:12mm) {$2$}; 
      \state                     (3) at (234:12mm) {$3$}; 
      \state                     (4) at (306:12mm) {$4$}; 
      \state                     (5) at (378:12mm) {$5$}; 

      \path (1) edge node {$a$} (2);
      \path (2) edge node {$a$} (3);
      \path (3) edge node {$a$} (4);
      \path (4) edge node {$a$} (5);
      \path (5) edge node {$a$} (1);
    \end{tikzpicture}
    &
    \begin{tikzpicture}[automaton, size=13mm, baseline=(2.base)]
      \state[initial, accepting] (1) at (0,.5) {$1$}; 
      \state                     (2) at (1,1)  {$2$}; 
      \state                     (3) at (1,0)  {$3$}; 
      \state[accepting]          (4) at (2,1)  {$4$}; 
      \state[accepting]          (5) at (2,0)  {$5$}; 
      
      \path (2) edge[loop above] node       {$b$} (2);
      \path (5) edge[loop below] node       {$a$} (5);
      \path (1) edge             node       {$a$} (2);
      \path (1) edge             node[swap] {$b$} (3);
      \path (2) edge             node       {$a$} (4);
      \path (3) edge             node       {$a$} (5);
      \path (3) edge             node       {$b$} (2);
    \end{tikzpicture}
  \end{tabular}
  
\end{exercice}

\begin{correction}

  \begin{question}
  \item 
    \begin{itemize}
    \item $\mathcal{S}(a)$ reconnu par
      \scalebox{.75}{\begin{tikzpicture}[shorten >=1pt,node distance=1.7cm,on grid,auto]
          \node[state, initial, initial text=] (0)              {}; 
          \node[state, accepting]              (1) [right=of 0] {}; 
          \path[-latex] (0) edge node {$a$} (1);
      \end{tikzpicture}}
      et $\mathcal{S}(b)$ reconnu par
      \scalebox{.75}{\begin{tikzpicture}[shorten >=1pt,node distance=1.7cm,on grid,auto]
          \node[state, initial, initial text=] (0)              {}; 
          \node[state, accepting]              (1) [right=of 0] {}; 
          \path[-latex] (0) edge node {$b$} (1);
      \end{tikzpicture}}
    \item $\mathcal{S}(a | b)$ reconnu par\\
      \scalebox{.75}{\begin{tikzpicture}[shorten >=1pt,node distance=1.7cm,on grid,auto]
          \node[state, initial, initial text=] (0)              {}; 
          \node[state]              (1) [above right=of 0] {}; 
          \node[state]              (2) [right=of 1] {}; 
          \node[state]              (3) [below right=of 0] {}; 
          \node[state]              (4) [right=of 3] {}; 
          \node[state, accepting]   (5) [above right=of 4] {}; 
          \path[-latex] (0) edge node {$\varepsilon$} (1);
          \path[-latex] (0) edge node {$\varepsilon$} (3);
          \path[-latex] (1) edge node {$a$} (2);
          \path[-latex] (3) edge node {$b$} (4);
          \path[-latex] (2) edge node {$\varepsilon$} (5);
          \path[-latex] (4) edge node {$\varepsilon$} (5);
      \end{tikzpicture}}
    \item $\mathcal{S}((a|b)^\star)$ reconnu par\\
      \scalebox{.75}{\begin{tikzpicture}[shorten >=1pt,node distance=1.7cm,on grid,auto]
          \node[state]                         (0)              {}; 
          \node[state]                         (1) [above right=of 0] {}; 
          \node[state]                         (2) [right=of 1] {}; 
          \node[state]                         (3) [below right=of 0] {}; 
          \node[state]                         (4) [right=of 3] {}; 
          \node[state]                         (5) [above right=of 4] {}; 
          \node[state, initial, initial text=] (6) [above left=of 0] {}; 
          \node[state, accepting]              (7) [below left=of 0] {}; 
          \path[-latex] (0) edge node {$\varepsilon$} (1);
          \path[-latex] (0) edge node {$\varepsilon$} (3);
          \path[-latex] (1) edge node {$a$} (2);
          \path[-latex] (3) edge node {$b$} (4);
          \path[-latex] (2) edge node {$\varepsilon$} (5);
          \path[-latex] (4) edge node {$\varepsilon$} (5);
          \path[-latex] (5) edge node {$\varepsilon$} (0);
          \path[-latex] (6) edge node {$\varepsilon$} (0);
          \path[-latex] (0) edge node {$\varepsilon$} (7);
      \end{tikzpicture}}
    \item $\mathcal{S}(L_1)$ reconnu par\\
      \scalebox{.75}{\begin{tikzpicture}[shorten >=1pt,node distance=1.7cm,on grid,auto]
          \node[state, initial, initial text=] (0) {}; 
          \node[state]                         (1) [right=of 0] {}; 
          \node[state]                         (2) [right=of 1] {}; 
          \node[state]                         (3) [right=of 2] {}; 
          \node[state]                         (4) [right=of 3] {}; 
          \node[state]                         (5) [right=of 4] {}; 
          \node[state]                         (6) [right=of 5] {}; 
          \node[state]                         (7) [right=of 6] {}; 
          \node[state]                         (8) [right=of 7] {}; 
          \node[state, accepting]              (9) [right=of 8] {};

          \node[state]                         (10) [below left=of 1]   {}; 
          \node[state]                         (11) [below right=of 1]  {}; 
          \node[state]                         (12) [below=of 10]       {}; 
          \node[state]                         (13) [below=of 11]       {}; 
          \node[state]                         (14) [below right=of 12] {}; 

          \node[state]                         (80) [below left=of 8]   {}; 
          \node[state]                         (81) [below right=of 8]  {}; 
          \node[state]                         (82) [below=of 80]       {}; 
          \node[state]                         (83) [below=of 81]       {}; 
          \node[state]                         (84) [below right=of 82] {}; 
          
          \path[-latex] (0) edge node {$\varepsilon$} (1);
          \path[-latex] (1) edge node {$\varepsilon$} (2);
          \path[-latex] (2) edge node {$\varepsilon$} (3);
          \path[-latex] (3) edge node {$a$} (4);
          \path[-latex] (4) edge node {$\varepsilon$} (5);
          \path[-latex] (5) edge node {$b$} (6);
          \path[-latex] (6) edge node {$\varepsilon$} (7);
          \path[-latex] (7) edge node {$\varepsilon$} (8);
          \path[-latex] (8) edge node {$\varepsilon$} (9);

          \path[-latex] (1) edge node {$\varepsilon$} (10);
          \path[-latex] (1) edge node {$\varepsilon$} (11);
          \path[-latex] (10) edge node {$a$} (12);
          \path[-latex] (11) edge node {$b$} (13);
          \path[-latex] (12) edge node {$\varepsilon$} (14);
          \path[-latex] (13) edge node {$\varepsilon$} (14);
          \path[-latex] (14) edge node {$\varepsilon$} (1);

          \path[-latex] (8) edge node {$\varepsilon$} (80);
          \path[-latex] (8) edge node {$\varepsilon$} (81);
          \path[-latex] (80) edge node {$a$} (82);
          \path[-latex] (81) edge node {$b$} (83);
          \path[-latex] (82) edge node {$\varepsilon$} (84);
          \path[-latex] (83) edge node {$\varepsilon$} (84);
          \path[-latex] (84) edge node {$\varepsilon$} (8);
      \end{tikzpicture}}
    \end{itemize}
  \item Même technique que pour $L_1$... 
  \item Il faut passer par le système d'équations rationnelles, puis utiliser le lemme d'Arden. \\
    On détaille ici la technique seulement pour l'automate en bas à gauche. 
    
    $\left\{\begin{array}{lllllllll}
    L_1 & = & a L_2 \mid b L_1 \mid\varepsilon &=& b^\star a L_2 \mid b^\star &=& b^\star a^+ b L_3 \mid b^\star\\
    L_2 & = & a L_2 \mid b L_3  &=& a^\star b L_3 \\
    L_3 & = & a L_4 \mid b L_1 \\
    L_4 & = & a L_2 \mid b L_3 \mid\varepsilon &&&=& a a^\star b L_3 \mid b L_3 \mid\varepsilon &=& a^\star b L_3 \mid \varepsilon\\
    \end{array}\right.$

    $L_3 = a L_4 \mid b L_1 = a^+ b L_3 \mid a \mid b^+ a^+ b L_3 \mid b^+ = b^\star a^+ b L_3 \mid a \mid b^+  = (b^\star a^+ b)^\star (a \mid b^+)$
    
    $\left\{\begin{array}{lllllllll}
    L_1 & = & b^\star a^+ b L_3 \mid b^\star     & = & (b^\star a^+ b)^+ (a \mid b^+) \mid b^\star \\     
    L_2 & = & a^\star b L_3                     & = & a^\star b (b^\star a^+ b)^\star (a \mid b^+)\\                       
    L_3 & = & (b^\star a^+ b)^\star (a \mid b^+) & = & (b^\star a^+ b)^\star (a \mid b^+) \\ 
    L_4 & = & a^\star b L_3 \mid \varepsilon    & = & a^\star b (b^\star a^+ b)^\star (a \mid b^+) \mid \varepsilon\\      
    \end{array}\right.$
    
  \item Réponses :
  \end{question}
  
  $\begin{array}[t]{cccc}
    aa\mid ba      &
    \varepsilon &
    (b\mid ab)^\star &
    (a\mid b)^\star ab (a\mid b)^\star\\
    (b^\star a^+ b)^+ (a \mid b^+) \mid b^\star&
    (a b^\star (a b^\star (a b^\star (a \mid b) \mid b) \mid b))^\star a b^\star a b^\star a&
    (aaaaa)^\star &
    ab^\star a\mid b a^+ \mid b b^+ a\\
  \end{array}$
  

\end{correction}

\endgroup
\endinput
